\chapter{Graphical editor, API and command line}\label{ch:interfaces}

\section{The graphical editor}
\code{python -m linkstudio serve} starts the web server and opens the editor in the browser.
\Cref{fig:gui-layout} shows its layout.

\begin{figure}[htbp]
\centering
\begin{tikzpicture}[x=1cm,y=1cm,font=\sffamily\scriptsize]
  \draw[black!40,fill=black!3,rounded corners=2pt] (0,0) rectangle (14,8.2);
  % top bar
  \draw[accent,fill=soft] (0.1,7.5) rectangle (13.9,8.1);
  \node[anchor=west] at (0.2,7.8) {\textbf{scenario name} \quad New \quad Save \quad Save as \quad Undo/Redo \quad Layout \quad ? \quad Export .py \quad seed};
  \node[anchor=east,accent3] at (13.8,7.8) {\textbf{valid}};
  % palette
  \draw[accent4,fill=accent4!6] (0.1,2.9) rectangle (2.6,7.4);
  \node[align=center] at (1.35,6.95) {\textbf{Blocks}\\search};
  \foreach \y/\t in {6.2/sources,5.75/encoders,5.3/mappers,4.85/modulators,4.4/media,3.95/decoders,3.5/measurements}
    \node[draw=black!25,fill=white,minimum width=2.1cm,minimum height=3.5mm] at (1.35,\y) {\t};
  % canvas
  \draw[black!30,fill=white] (2.7,2.9) rectangle (10.6,7.4);
  \node at (6.65,7.15) {\textbf{Canvas}};
  \foreach \x/\t/\c in {3.4/src/kbits,4.9/map/ksym,6.4/awgn/ksym,7.9/demap/kllr,9.6/ber/gray}
    \node[draw=\c,fill=\c!10,rounded corners=1pt,minimum width=1.1cm,minimum height=6mm] (n\t) at (\x,5.2) {\t};
  \draw[->,kbits,thick] (nsrc)--(nmap); \draw[->,ksym,thick] (nmap)--(nawgn);
  \draw[->,ksym,thick] (nawgn)--(ndemap); \draw[->,kllr,thick] (ndemap)--(nber);
  \draw[->,kbits,thick] (nsrc.south) -- ++(0,-0.6) -| (nber.south);
  \draw[black!30,fill=black!5] (9.2,3.0) rectangle (10.5,3.8); \node at (9.85,3.4) {minimap};
  % inspector
  \draw[accent2,fill=accent2!5] (10.7,2.9) rectangle (13.9,7.4);
  \node[align=left,anchor=north west] at (10.8,7.3) {\textbf{Inspector}\\[2pt] name\\ parameters (forms,\\ sliders, \texttt{auto})\\ frequency response\\[3pt] Flip \;\; Duplicate\\ Delete};
  % dock
  \draw[accent3,fill=accent3!6] (0.1,0.1) rectangle (13.9,2.8);
  \node[anchor=north west] at (0.2,2.75) {\textbf{Live} \quad Simulate \quad Sweep};
  \draw[black!30,fill=white] (0.3,0.3) rectangle (3.0,2.2); \node at (1.65,1.25) {metrics + CI};
  \foreach \x/\t in {3.2/constellation,5.9/spectrum,8.6/eye diagram,11.3/BER curve}
    { \draw[black!30,fill=white] (\x,0.3) rectangle ++(2.5,1.9); \node at (\x+1.25,1.25) {\t}; }
\end{tikzpicture}
\caption{Layout of the editor: palette, canvas, inspector and the run dock with its three
modes.}
\label{fig:gui-layout}
\end{figure}

\paragraph{Palette.}
The \numblocks{} blocks are grouped by functional class in signal-chain order: sources, encoders, bit
processing, mappers, multiplexers, framing, modulators, filters and resamplers, amplifiers and
attenuators, transmission media, impairments, then the receiver classes (demodulators,
synchronization, equalizers, detectors and estimators, demultiplexers, demappers, decoders) and
the measurements. Complementary classes (encoders and decoders, mappers and demappers, \dots)
share a colour. Only the large classes are split into subclasses (code families, media); the
others are flat lists. Every block draws its function on the canvas: codewords with parity, the
shift register and trellis of convolutional codes, the Tanner graph, the interleaver permutation
drawn from its seed (or the block matrix numbered in write order), the scrambler LFSR with its
real taps and seed, pulse shapes, constellations, and the physical blocks (adders, mixers,
antennas, fibers, lasers) as drawings without a frame.

\paragraph{Editing.}
\begin{itemize}
  \item \textbf{Adding blocks:} drag a block from the palette or double-click it.
  \item \textbf{Connecting:} drag from an output handle to an input handle. Handle colours are
  the signal kinds of \cref{tab:kinds}; incompatible connections are refused, and dashed
  handles are optional inputs.
  \item \textbf{Selection:} click a block; Shift or Ctrl + click adds blocks, Shift + drag on
  the empty canvas selects the blocks in a box, Ctrl+A selects all and Esc clears the
  selection. With several blocks selected a bar over the canvas offers \emph{Align row} (the
  links between them become straight), \emph{Align column}, \emph{Distribute} (equal gaps),
  \emph{Flip}, \emph{Duplicate} and \emph{Delete}.
  \item \textbf{Clipboard:} Ctrl+C, Ctrl+X and Ctrl+V copy, cut and paste blocks with the
  connections between them; the copy is placed at the mouse and pasting again steps the
  copies. Ctrl+D duplicates, Delete deletes, the arrow keys move the selection (Shift + arrow
  by a larger step).
  \item \textbf{Direction:} a block faces right (inputs on the left, outputs on the right) or
  left. F, or \emph{Flip} in the inspector, turns the selected blocks around. A block dragged
  to the other side of the blocks it is linked to turns around by itself, so a chain can run
  right to left or turn back on itself without moving any port. The block drawings are
  mirrored with the block; their text stays readable and plots against time or frequency keep
  their axes.
  \item \textbf{Moving ports:} Ctrl + drag a handle to move it anywhere on the block border,
  e.g.\ for feedback paths. Port places and directions are saved in the scenario
  (\code{port_layout}).
  \item \textbf{Alignment:} while a block is dragged it snaps so that its connected ports line
  up with the ports at the other end (straight links), or else so that its edges or centre line
  up with a neighbouring block; a guide line shows the alignment.
  \item \textbf{Automatic layout} (Ctrl+L) arranges the blocks along the signal flow
  (\cref{fig:auto-layout}). Each block goes in the column given by its longest distance from a
  source, and sources that feed late blocks are pulled forward. The longest path (the main
  chain) stays on one straight line. Branches such as training sequences, channel estimates
  and meters go in lanes under it, ordered by barycentre sweeps to reduce crossings. A long
  chain wraps into rows sized from the shape of the canvas and run alternately left to right
  and right to left (a horseshoe); the blocks of the reversed rows face left. Scopes and meters
  that tap a link from their bottom side sit above the gap after the block they tap.
  Unconnected blocks go in a last row.
  \item \textbf{Links} are drawn with orthogonal segments routed around the blocks (A$^\ast$
  search on a sparse grid of block edges and channels, with a penalty per bend), so they never
  pass behind a block. All links are routed together. Where links of different nets would run
  on top of each other along a channel, their shared segments are moved apart so that each
  link stays visible.
  \item \textbf{Inspector:} parameter forms are generated from each block's schema. Bounded and
  dB-valued parameters get sliders, and \code{auto} leaves a value to upstream metadata. For
  linear blocks the inspector also plots the \emph{frequency response} (magnitude, phase or group
  delay): the model transfer function of cables, CTLE, lasers and VCSELs, photodiodes, single-
  and multimode fibers and pulse shapers, and an estimate $\hat H=S_{xy}/S_{xx}$ of one
  realization for random multipath channels.
  \item \textbf{Undo:} Ctrl+Z and Ctrl+Y undo and redo any change: blocks, connections,
  parameters, positions, port places and directions, scopes.
  \item \textbf{Shortcuts:} ? (or the ? button) lists them all. Besides those above: Ctrl+S
  saves, Ctrl+Shift+S saves as, Ctrl+O opens the scenario browser, Ctrl+B and Ctrl+I show or hide
  the block palette and the inspector.
  Saving keeps the run settings (\code{simulation}, \code{sweep}) and the ports of a subsystem
  file, so a scenario saved from the GUI runs the same from the command line.
\end{itemize}

\begin{figure}[htbp]
\centering
\begin{tikzpicture}[x=1cm,y=1cm,font=\sffamily\scriptsize,
    b/.style={draw=accent,fill=soft,rounded corners=1pt,minimum width=1.25cm,minimum height=6mm},
    l/.style={->,thick,black!60}]
  \foreach \i/\t in {0/src,1/enc,2/map,3/mod,4/tx}
    \node[b] (r\i) at (1.9*\i,2) {\t};
  \foreach \i/\t in {4/fiber,3/rx,2/eq,1/demap,0/dec}
    \node[b] (q\i) at (1.9*\i,0) {\t};
  \node[b,draw=gray,fill=gray!8] (ber) at (-1.9,0) {ber};
  \node[b,draw=gray,fill=gray!8,minimum width=1cm] (sc) at (2.85,2.9) {scope};
  \foreach \i/\j in {0/1,1/2,2/3,3/4} \draw[l] (r\i) -- (r\j);
  \draw[l] (r4.east) -- ++(0.35,0) |- (q4.east);
  \foreach \i/\j in {4/3,3/2,2/1,1/0} \draw[l] (q\i) -- (q\j);
  \draw[l] (q0) -- (ber);
  \draw[l] (r0.south) -- ++(0,-0.4) -| (ber.north);
  \fill[black!60] (2.85,2) circle (1.3pt);
  \draw[thick,black!60] (2.85,2) -- (sc.south);
  \node[anchor=west,commentcolor] at (8.9,2) {row 1: faces right};
  \node[anchor=west,commentcolor] at (8.9,0) {row 2: faces left};
\end{tikzpicture}
\caption{Automatic layout of a long chain: it wraps into a second row that runs right to left
with its blocks facing left. A scope tapping a link from its bottom side sits above it.}
\label{fig:auto-layout}
\end{figure}

The run buttons are disabled while the graph is invalid, and the validator's message is shown
next to them.

\paragraph{Modes.}
\begin{description}[font=\sffamily\bfseries]
  \item[Live] simulates one frame per step (batch size 1), 1--30 steps per second, and streams
  every probe; the metrics accumulate across steps. Parameter changes are applied between steps
  and the meters restart, so the effect of a slider is visible immediately. Structural edits
  restart the session. The rate only sets how often the scopes refresh, not the results.
  \item[Simulate] runs a Monte-Carlo job in the background, either for a fixed number of
  batches or until a metric reaches a threshold. \emph{frames/batch} sets the batch size (saved in
  the scenario) and \emph{cores} the number of worker processes (\cref{sec:parallel}). It shows the
  final metrics (BER with its 95\,\% confidence interval) and the probes of the last batch.
  \item[Sweep] runs one or two nested parameter sweeps and plots any metric, with an optional
  theoretical curve; with several cores the points run in parallel. Curves can be kept and
  overlaid, and the results exported as CSV or PNG.
\end{description}

\paragraph{The Scope block.}
Signals are displayed by the \blk{scope} block (class \emph{measurements}), an instrument with a
screen. Its inputs enter from the bottom side: drop the scope from the palette onto a link, or
drag one of its free inputs onto a link, to tap the output that feeds it. The tap is drawn as a
junction dot just after that output. A scope has up to four inputs; the selector in its corner
chooses the displayed one, and the bottom-right corner resizes it. The plot selector under its
name (parameter \prm{view}) lists the views that suit the displayed signal:
\begin{center}\small
\begin{tabular}{@{}ll@{}}
\toprule
Kind & Views (the first is the \code{auto} choice)\\
\midrule
\kind{SAMPLES} & spectrum (Welch), constellation, eye diagram, time\\
\kind{SYMBOLS} & constellation, time, spectrum\\
\kind{GRID} & constellation, resource-grid magnitude\\
\kind{CSI} & impulse / frequency response\\
\kind{LLR} & histogram, time\\
\kind{BITS} & time\\
\bottomrule
\end{tabular}
\end{center}
A view that does not suit the signal falls back to the first one of its kind. After a
\emph{Simulate} run every view is already computed, so switching needs no new run; in
\emph{Live} the next frame shows the new view, and the measurements are not restarted. Older
scenarios with one scope type per view (\code{scope_constellation}, \dots) are converted to
\blk{scope} with the matching view when they are loaded.

Axes carry physical units. The server sends values in base SI units together with their unit,
taken from the \meta{unit} metadata: $\sqrt{\mathrm{W}}$ for optical fields, A
for photocurrents, none for normalized baseband signals. The editor picks the SI prefix
(ps, GHz, mA, \dots) so that tick labels stay short. Spectra are absolute power spectral
densities (dBm/Hz for optical fields, one-sided for real signals). The eye diagram shows time
in seconds over two unit intervals and its vertical axis is centred on the mean of the traces
and spans all of them.

\section{The HTTP and WebSocket API}
The editor talks to the FastAPI application in \file{linkstudio/api/app.py}. Everything it does
can also be scripted.

\begin{center}\small
\begin{tabularx}{\textwidth}{@{}llX@{}}
\toprule
Method & Path & Purpose\\
\midrule
GET  & \code{/api/health} & server version and number of CPU cores\\
GET  & \code{/api/blocks} & block catalogue: documentation, ports, parameter JSON Schema\\
GET  & \code{/api/scenarios} & scenario list with name, category, description, block count\\
GET/PUT & \code{/api/scenarios/{path}} & load / save a scenario in editor format\\
POST & \code{/api/validate} & validate an editor graph; returns errors and port kinds\\
POST & \code{/api/layout} & server-side layout (the editor lays out locally)\\
POST & \code{/api/export/python} & standalone Python script equivalent to the graph\\
POST & \code{/api/response} & frequency response of one block (magnitude, phase, group delay)\\
POST & \code{/api/jobs} & start a simulate or sweep job\\
GET  & \code{/api/jobs/{id}} & job status, progress, partial and final results\\
POST & \code{/api/jobs/{id}/cancel} & cancel a job\\
WS   & \code{/api/live} & interactive live mode\\
\bottomrule
\end{tabularx}
\end{center}

Jobs run in a thread pool (\file{api/jobs.py}); the \code{workers} entry of the job
configuration sends the batches or sweep points to that many worker processes. Their results
contain JSON-safe metrics (NaN
becomes \code{null}), the probe data of the last batch and, for sweeps of a single parameter
whose value matches a known modulation, theoretical BER curves. A failed job reports the
message \code{failed} with the details in a separate \code{detail} field.

\begin{figure}[htbp]
\centering
\begin{tikzpicture}[x=1cm,y=-0.62cm,font=\sffamily\scriptsize]
  \node[blockbox] (c) at (0,0) {browser};
  \node[blockbox] (s) at (7,0) {server (\code{/api/live})};
  \draw[black!40] (0,0.6) -- (0,10.5); \draw[black!40] (7,0.6) -- (7,10.5);
  \draw[->] (0,1.5) -- node[above]{\code{start} \{graph, rate\}} (7,1.5);
  \draw[<-] (0,2.5) -- node[above]{\code{started}} (7,2.5);
  \draw[<-,accent3] (0,3.6) -- node[above]{\code{frame} \{batch, metrics, probes, step\_ms\}} (7,3.6);
  \draw[<-,accent3] (0,4.6) -- node[above]{\code{frame}} (7,4.6);
  \draw[->,accent2] (0,5.6) -- node[above]{\code{set_param} \{path, value\}} (7,5.6);
  \node[anchor=west,note] at (7.1,6.3) {applied before the next batch;\\meters restart};
  \draw[<-,accent3] (0,7.0) -- node[above]{\code{frame}} (7,7.0);
  \draw[->] (0,8.0) -- node[above]{\code{probes} / \code{rate} / \code{reset}} (7,8.0);
  \draw[->] (0,9.2) -- node[above]{\code{stop}} (7,9.2);
  \draw[<-] (0,10.2) -- node[above]{\code{stopped}} (7,10.2);
\end{tikzpicture}
\caption{Message flow of the live mode. The server steps the simulation in a worker thread and
paces frames to the requested rate.}
\end{figure}

\section{Command line}
\begin{center}\small
\begin{tabularx}{\textwidth}{@{}lX@{}}
\toprule
\code{python -m linkstudio blocks [--category C]} & list block types\\
\code{python -m linkstudio describe TYPE} & ports, parameters and documentation of a block\\
\code{python -m linkstudio run FILE.yaml [--csv OUT]} & run a scenario (simulation or sweep)\\
\code{python -m linkstudio serve [--host H] [--port P] [--no-browser]} & start the editor\\
\code{python scripts/main.py [--quick]} & run all demos with theory checks\\
\bottomrule
\end{tabularx}
\end{center}
The demo scripts in \file{scripts/} (one per subsystem) are short, readable examples of the
Python API. Each prints simulated values next to theory and returns a pass/fail status.
